% ==============================================================================
% MÓDULO CAPÍTULO II: ENCABEZADO Y AMBIENTE NATURAL FENOMENOLÓGICO
% Archivo: cap_2/2.0_encabezado.tex
% Enfoque: 100% Calidad de Atención (Donabedian y Watson)
% ==============================================================================
\section{MARCO CONTEXTUAL}
\label{sec:cap2_contexto}

En la investigación cualitativa bajo el diseño fenomenológico, el marco contextual no se reduce a una enumeración estática de datos cuantitativos ni a un inventario geográfico desvinculado; constituye la reconstrucción densa, profunda y holística del \textit{ambiente natural} donde transcurre la vida cotidiana de las personas y donde se gesta el fenómeno estudiado \cite{sampieri2018, creswell2013}. Para comprender la esencia de la experiencia vivida (\textit{Lebenswelt}) de los usuarios respecto a la calidad de atención y el cuidado humano recibido en el Centro de Salud Belén, es indispensable caracterizar la atmósfera física, cultural, social y sanitaria que envuelve, condiciona y dota de sentido a sus sentires, expectativas, frustraciones y valoraciones intersubjetivas en el encuentro asistencial \cite{vanmanen1990, diresa2025, donabedian2005, watson2008}.
